\documentclass[11pt,a4paper]{amsart}
\pdfinfoomitdate=1
\pdftrailerid{}
\pdfsuppressptexinfo=15
\usepackage[T1]{fontenc}
\usepackage{lmodern}
\usepackage[margin=27mm]{geometry}
\usepackage{microtype}
\usepackage{amsmath,amssymb,mathtools}
\usepackage{booktabs}
\usepackage{xcolor}
\usepackage[numbers,sort&compress]{natbib}
\usepackage[colorlinks=true,linkcolor=blue!55!black,citecolor=blue!55!black,urlcolor=blue!55!black]{hyperref}
\usepackage[nameinlink,noabbrev]{cleveref}
\raggedbottom

\newtheorem{theorem}{Theorem}[section]
\newtheorem{proposition}[theorem]{Proposition}
\newtheorem{corollary}[theorem]{Corollary}
\newtheorem{lemma}[theorem]{Lemma}
\theoremstyle{definition}
\newtheorem{definition}[theorem]{Definition}
\theoremstyle{remark}
\newtheorem{remark}[theorem]{Remark}

\newcommand{\Z}{\mathbb Z}
\newcommand{\cL}{\mathcal L}

\title[Critical-Jordan rectangle discrepancy]{Exact Dyadic Rectangle Discrepancy at a Critical Jordan Resonance}
\author{Anonymous}
\date{Internal Stage 2 draft, 27 August 2026}
\subjclass[2020]{37B10, 37A30, 68R15}
\keywords{multidimensional substitution, discrepancy, Jordan block, recognizability, dyadic rectangle, symbolic dynamics}

\begin{document}

\begin{abstract}
We give an explicit primitive, aperiodic $2\times2$ block substitution whose
incidence matrix has a size-two Jordan block at the critical eigenvalue $2$,
equal to the linear expansion.  A radius-one finite certificate recognizes
the complete supertile grid.  For the centered observable
$f=\mathbf1_A-\mathbf1_B$, every translated dyadic rectangle of side lengths
$2^m$ and $2^n$ satisfies a sharp discrepancy law:
\[
 \sup_{\substack{x\in X_\vartheta\\t\in\Z^2}}
 \left|\sum_{q\in t+[0,2^m)\times[0,2^n)}f(x_q)\right|
 =2^{\max(m,n)}\left(1+\frac{\min(m,n)}2\right).
\]
The proof combines an exact generalized-eigenvector identity with a
pointwise periodic envelope.  Thus the critical Jordan logarithm is not only
visible on complete supertiles: its sharp worst-case coefficient is determined
uniformly over translations and for every dyadic aspect ratio.
\end{abstract}

\maketitle

\section{Introduction}

For primitive substitution systems, incidence eigenvalues govern deviations
of letter counts from their invariant frequencies.  This principle is
classical in one-dimensional symbolic discrepancy
\cite{Adamczewski2004} and has far-reaching higher-dimensional tiling forms
\cite{BufetovSolomyak2013}.  At a boundary eigenvalue---modulus equal to the
linear boundary scaling---a nontrivial Jordan block can introduce logarithmic
losses.  General theory gives robust bounds, but it need not determine the
sharp coefficient for a concrete symbolic window.

We exhibit a three-letter planar block substitution for which the critical
mechanism can be computed exactly.  Two finite structures cooperate.  First,
the incidence matrix has a length-two left Jordan chain at eigenvalue $2$.
Second, the three level-$n$ observable arrays admit pointwise lower and upper
envelopes that are $2^n$-periodic.  A translated $2^n$ square samples each
residue once, turning those pointwise envelopes into sharp sum bounds.

The recognizability assertion is also explicit.  The $20$ legal $2\times2$
patches form a substitution-closed set, while the $63$ legal $3\times3$
patches split into four disjoint phase classes.  This is a complete
radius-one certificate, rather than an appeal to recognizability in the
abstract.  General recognizability and its computability for
multidimensional constant-shape substitutions have been developed much more
broadly \cite{Solomyak1998,Cabezas2023,CabezasLeroy2025}; those general results
are not claimed here.

\section{The substitution and its subshift}

Let $\vartheta$ be the $2\times2$ substitution on $\{A,B,C\}$ given by
\begin{equation}\label{eq:substitution}
 \vartheta(A)=\begin{pmatrix}A&A\\B&C\end{pmatrix},\qquad
 \vartheta(B)=\begin{pmatrix}B&B\\B&C\end{pmatrix},\qquad
 \vartheta(C)=\begin{pmatrix}A&C\\C&C\end{pmatrix}.
\end{equation}
Let $X_\vartheta\subset\{A,B,C\}^{\Z^2}$ be the subshift whose finite patches
occur in iterates $\vartheta^n(a)$.  Rows index output letters and columns
index parent letters in the incidence matrix
\begin{equation}\label{eq:matrix}
 M=\begin{pmatrix}
 2&0&1\\
 1&3&0\\
 1&1&3
 \end{pmatrix}.
\end{equation}
Since
\[
 M^2=\begin{pmatrix}5&1&5\\5&9&1\\6&6&10\end{pmatrix}>0,
\]
the substitution is primitive.  Its Perron eigenvalue is $4$, and its
normalized right Perron vector is
\begin{equation}\label{eq:frequencies}
 \nu=(1/4,1/4,1/2)^{\mathsf T}.
\end{equation}

\subsection{A finite recognizability certificate}

Write $\cL_2$ for the set of legal $2\times2$ patches.  Starting with the
three images in \eqref{eq:substitution}, repeatedly adjoin every $2\times2$
subpatch of $\vartheta(P)$ for every current patch $P$.  The cardinalities are
\[
 3\longrightarrow13\longrightarrow20\longrightarrow20.
\]
Thus the resulting $20$-patch set is substitution-closed.  Every fine
$2\times2$ patch crosses at most a $2\times2$ parent patch, so induction shows
that this fixed set is exactly $\cL_2$.

Every legal $3\times3$ patch is contained in $\vartheta(P)$ for some
$P\in\cL_2$.  Classify it by the location of its upper-left corner modulo
$2\Z^2$.  With one choice of origin, the class sizes in phase order
$(0,0),(0,1),(1,0),(1,1)$ are
\begin{equation}\label{eq:phase-counts}
 18,\quad20,\quad15,\quad10,
\end{equation}
and the four classes are pairwise disjoint.  The accompanying script performs
this exhaustive set calculation with exact strings.

\begin{proposition}[radius-one recognizability]\label{prop:recognizable}
Every $3\times3$ patch in $X_\vartheta$ determines its phase modulo $2\Z^2$.
Consequently every $x\in X_\vartheta$ has a unique level-one phase--parent
pair: a unique offset $\tau\in\{0,1\}^2$ for the supertile grid and a unique
parent configuration.  Iterating gives a unique level-$n$ cutting for every
$n$.
\end{proposition}

\begin{proof}
Take substituted approximants containing larger and larger centered patches
of $x$.  Pass to a subsequence with a fixed grid offset and then take a
compact limit of their parent configurations; this gives a level-one cutting.
If its offset is $\tau$, a $3\times3$ window with upper-left corner $z$
belongs to the finite phase class indexed by
$(z-\tau)\bmod2\Z^2$.  For a second cutting with offset $\tau'$, the same
window would belong to the class indexed by
$(z-\tau')\bmod2\Z^2$.  Pairwise disjointness of the four classes forces
$\tau=\tau'$.  At each common anchor, the corresponding $2\times2$ block is
one of the three distinct images in \eqref{eq:substitution}, so it uniquely
determines the parent letter.  This proves uniqueness at level one, and
iteration proves the general case.
\end{proof}

\begin{corollary}\label{cor:aperiodic}
The subshift $X_\vartheta$ is aperiodic, minimal, and uniquely ergodic, with
letter frequencies \eqref{eq:frequencies}.
\end{corollary}

\begin{proof}
If $p\in\Z^2$ is a period, the intrinsic phase detector forces
$p\in2\Z^2$.  Desubstitution then makes $p/2$ a period.  Iteration forces $p$
to be divisible by $2^n$ for every $n$, hence $p=0$.  Primitivity gives
minimality and uniform patch frequencies, and therefore unique ergodicity;
the letter frequencies are the normalized Perron vector.
\end{proof}

\section{The critical Jordan chain}

Let
\[
 f(A)=1,\qquad f(B)=-1,\qquad f(C)=0.
\]
By \eqref{eq:frequencies}, $f$ has invariant mean zero.  Define the left row
vectors
\[
 w_1=(1,-1,0),\qquad w_0=(-1,-1,1).
\]
Direct multiplication gives
\begin{equation}\label{eq:jordan}
 w_0M=2w_0,\qquad w_1M=2w_1+w_0,
\end{equation}
while
\[
 \det(\lambda I-M)=(\lambda-4)(\lambda-2)^2.
\]
Thus eigenvalue $2$ has a size-two Jordan block.  From
\eqref{eq:jordan},
\begin{equation}\label{eq:power}
 w_1M^n=2^nw_1+n2^{n-1}w_0.
\end{equation}

For $a\in\{A,B,C\}$ let $S_n(a)$ be the sum of $f$ over the complete
$n$-supertile $\vartheta^n(a)$.  Evaluating \eqref{eq:power} on the three
coordinate vectors gives the exact identities
\begin{equation}\label{eq:supertile-sums}
\begin{aligned}
 S_n(A)&=2^n\left(1-\frac n2\right),\\
 S_n(B)&=-2^n\left(1+\frac n2\right),\\
 S_n(C)&=n2^{n-1}.
\end{aligned}
\end{equation}
The eigenvalue modulus $2$ equals the linear expansion, hence the designation
\emph{critical}.  The Jordan factor $n$ becomes $\log_2R$ at side length
$R=2^n$.

\section{A periodic-envelope lemma}

Let $F_n^a$ be the $2^n\times2^n$ array obtained by applying $f$ entrywise
to $\vartheta^n(a)$.  The block recursions from \eqref{eq:substitution} are
\[
 F_{n+1}^A=\begin{pmatrix}F_n^A&F_n^A\\F_n^B&F_n^C\end{pmatrix},\quad
 F_{n+1}^B=\begin{pmatrix}F_n^B&F_n^B\\F_n^B&F_n^C\end{pmatrix},\quad
 F_{n+1}^C=\begin{pmatrix}F_n^A&F_n^C\\F_n^C&F_n^C\end{pmatrix}.
\]

\begin{lemma}[pointwise envelope]\label{lem:envelope}
For every $n\geq0$,
\[
 F_n^B\leq F_n^A,\qquad F_n^B\leq F_n^C
\]
entrywise.  Moreover $F_n^A\leq F_n^C$ except at the single position
$p_n=(0,2^n-1)$, where $F_n^A(p_n)-F_n^C(p_n)=1$.
\end{lemma}

\begin{proof}
The assertion is immediate at $n=0$.  Substitute the three displayed block
recursions.  Each required comparison reduces to the induction hypothesis in
one of the four blocks.  The unique exceptional position moves into the
upper-right block, from $p_n$ to $p_{n+1}=(0,2^{n+1}-1)$.
\end{proof}

Extend $F_n^B$ periodically to $\Z^2$, and likewise extend
\[
 U_n=F_n^C+\mathbf1_{\{p_n\}}.
\]
By unique level-$n$ cutting and \cref{lem:envelope}, after shifting their
common phase these arrays bound the entire observable field of every
$x\in X_\vartheta$ pointwise.  Every square
$Q=t+[0,2^n)^2$, $t\in\Z^2$, samples each residue class modulo
$2^n\Z^2$ exactly once.  Therefore
\begin{equation}\label{eq:square-bounds}
 -2^n\left(1+\frac n2\right)
 \leq \sum_{q\in Q}f(x_q)
 \leq n2^{n-1}+1.
\end{equation}
The left endpoint is attained by a complete $B$-supertile.

\section{Exact discrepancy on every dyadic rectangle}

For $m,n\geq0$, define
\[
 D(m,n)=\sup\left\{
 \left|\sum_{q\in t+[0,2^m)\times[0,2^n)}f(x_q)\right|:
 x\in X_\vartheta,\ t\in\Z^2
 \right\}.
\]

\begin{theorem}[sharp dyadic rectangle law]\label{thm:rectangle}
For all $m,n\geq0$,
\begin{equation}\label{eq:rectangle}
 D(m,n)=2^{\max(m,n)}
 \left(1+\frac{\min(m,n)}2\right).
\end{equation}
\end{theorem}

\begin{proof}
First take $m=n=k$.  The absolute value of the lower endpoint in
\eqref{eq:square-bounds} dominates the upper endpoint, and a complete
$B$-supertile attains it.  Hence
$D(k,k)=2^k(1+k/2)$.

Suppose $m=k\leq n=\ell$.  Partition a
$2^k\times2^\ell$ rectangle into $2^{\ell-k}$ dyadic squares of side $2^k$.
The triangle inequality and the square case give
\[
 D(k,\ell)\leq2^{\ell-k}D(k,k)=2^\ell(1+k/2).
\]
For equality, take a complete $\ell$-supertile of type $B$ and its top strip
of height $2^k$.  The top row of the parent array
$\vartheta^{\ell-k}(B)$ consists entirely of $B$'s, so the strip is a row of
$2^{\ell-k}$ complete $k$-supertiles of type $B$ and attains the negative
bound.  If $n=k\leq m=\ell$, use the left strip instead; the left column of
$\vartheta^{\ell-k}(B)$ also consists entirely of $B$'s.
\end{proof}

In particular, square windows of side $R=2^n$ have exact worst-case
discrepancy
\[
 D(n,n)=R\left(1+\frac12\log_2R\right).
\]
The logarithm is therefore a global translated-window phenomenon, not merely
an ancestral-supertile calculation.

\section{Scope and reproducibility}

The finite control checks primitivity, the Jordan identities, the supertile
formulas through level $32$, and the pointwise envelope through level $8$.
Those level-bounded tests are regression checks for the displayed algebra and
induction.  Separately, starting from the three substitution images, the
script constructs the complete $2\times2$ patch closure
$3\to13\to20\to20$ and then all $63$ legal $3\times3$ patches, verifying that
their four phase classes are pairwise disjoint.  This patch-language
calculation is the exhaustive finite base of the recognizability proof.

The result should not be read as a general discrepancy theorem.  General
spectral control of substitution discrepancy belongs to the existing theory
\cite{Adamczewski2004,BufetovSolomyak2013}.  In the planar critical
size-two-Jordan regime, the general Lipschitz-domain estimate of
\citet[Corollary~4.5]{BufetovSolomyak2013} allows
$O(R(\log R)^2)$; the special pointwise order in
\eqref{eq:substitution} is what yields the exact $R\log R$ law here.  We claim
only this explicit system, its radius-one certificate, its exact Jordan
chain, and formula \eqref{eq:rectangle}.  External priority and release
questions remain outside this internal draft.

\bibliographystyle{plainnat}
\bibliography{references}

\end{document}
